\documentclass[11pt,fontset=fandol]{ctexart}
\usepackage[a4paper,margin=23mm]{geometry}
\usepackage{amsmath,amssymb,mathtools,bm,mathrsfs,esint}
\usepackage{graphicx,longtable,booktabs,array,calc}
\usepackage[normalem]{ulem}
\usepackage{xcolor,hyperref}
\hypersetup{hidelinks}
\setcounter{secnumdepth}{-1}
\newcounter{none}
\providecommand{\tightlist}{\setlength{\itemsep}{0pt}\setlength{\parskip}{0pt}}
\providecommand{\passthrough}[1]{#1}
\usepackage{amssymb,mathrsfs,esint,bm,newunicodechar}
\allowdisplaybreaks
\setlength{\emergencystretch}{3em}
\providecommand{\llangle}{\langle\!\langle}
\providecommand{\rrangle}{\rangle\!\rangle}
\newunicodechar{∞}{\ensuremath{\infty}}
\newunicodechar{∈}{\ensuremath{\in}}
\newunicodechar{∑}{\ensuremath{\sum}}
\newunicodechar{∫}{\ensuremath{\int}}
\newunicodechar{≤}{\ensuremath{\leq}}
\newunicodechar{≥}{\ensuremath{\geq}}
\newunicodechar{≠}{\ensuremath{\neq}}
\newunicodechar{→}{\ensuremath{\to}}
\newunicodechar{⇒}{\ensuremath{\Rightarrow}}
\newunicodechar{α}{\ensuremath{\alpha}}
\newunicodechar{β}{\ensuremath{\beta}}
\newunicodechar{γ}{\ensuremath{\gamma}}
\newunicodechar{λ}{\ensuremath{\lambda}}
\newunicodechar{φ}{\ensuremath{\varphi}}
\newunicodechar{ε}{\ensuremath{\varepsilon}}
\newunicodechar{σ}{\ensuremath{\sigma}}
\newunicodechar{θ}{\ensuremath{\theta}}
\newunicodechar{δ}{\ensuremath{\delta}}
\newunicodechar{η}{\ensuremath{\eta}}
\newunicodechar{π}{\ensuremath{\pi}}
\newunicodechar{ξ}{\ensuremath{\xi}}
\newunicodechar{Σ}{\ensuremath{\Sigma}}
\newunicodechar{Δ}{\ensuremath{\Delta}}
\newunicodechar{∂}{\ensuremath{\partial}}
\newunicodechar{∀}{\ensuremath{\forall}}
\newunicodechar{∃}{\ensuremath{\exists}}
\newunicodechar{√}{\ensuremath{\surd}}
\newunicodechar{−}{\ensuremath{-}}
\newunicodechar{μ}{\ensuremath{\mu}}
\newunicodechar{ρ}{\ensuremath{\rho}}
\newunicodechar{τ}{\ensuremath{\tau}}
\newunicodechar{ω}{\ensuremath{\omega}}
\newunicodechar{Ω}{\ensuremath{\Omega}}
\newunicodechar{χ}{\ensuremath{\chi}}
\newunicodechar{ψ}{\ensuremath{\psi}}
\newunicodechar{Γ}{\ensuremath{\Gamma}}
\newunicodechar{Φ}{\ensuremath{\Phi}}
\newunicodechar{Λ}{\ensuremath{\Lambda}}
\newunicodechar{ℝ}{\ensuremath{\mathbb{R}}}
\newunicodechar{ℤ}{\ensuremath{\mathbb{Z}}}
\newunicodechar{ℕ}{\ensuremath{\mathbb{N}}}
\newunicodechar{△}{\ensuremath{\triangle}}
\newunicodechar{⊗}{\ensuremath{\otimes}}
\begin{document}
\section{2021年杭州电子科技大学601数学分析考研真题}\label{ux5e74ux676dux5ddeux7535ux5b50ux79d1ux6280ux5927ux5b66601ux6570ux5b66ux5206ux6790ux8003ux7814ux771fux9898}

来源：\texttt{HDU/601数学分析/2021/2021年杭州电子科技大学601数学分析考研真题.pdf}。原
PDF 页序：1, 2。

\section{杭州电子科技大学}\label{ux676dux5ddeux7535ux5b50ux79d1ux6280ux5927ux5b66}

\section{2021年攻读硕士学位研究生招生考试}\label{ux5e74ux653bux8bfbux7855ux58ebux5b66ux4f4dux7814ux7a76ux751fux62dbux751fux8003ux8bd5}

\section{《数学分析》试题}\label{ux6570ux5b66ux5206ux6790ux8bd5ux9898}

(试题共四大题，共2页，总分150分）

姓名 报考专业 准考证号

【所有答案必须写在答题纸上，做在试卷或草稿纸上无效！】

\subsection{一、填空题（本大题共8小题，每小题5分，共40分)}\label{ux4e00ux586bux7a7aux9898ux672cux5927ux9898ux51718ux5c0fux9898ux6bcfux5c0fux98985ux5206ux517140ux5206}

\begin{enumerate}
\def\labelenumi{\arabic{enumi}.}
\tightlist
\item
  \(\lim _ { n  \infty } \frac { n ^ { 2 } } { 2 ^ { n } } = (\) ）.
\end{enumerate}

2.设点(0,1)是曲线 \(y = x ^ { 3 } + b x ^ { 2 } + 1\) 的拐点，则b=（ ）.

3.设
\(\left\{ \begin{array} { l } { x = a \cos ^ { 3 } t } \\ { y = a \sin ^ { 3 } t , } \end{array} \right.\)
则
\(\left. { \frac { d y } { d x } } \right| _ { x = \frac { \pi } { 4 } } = (\)
）.

4.曲线y=x(x-1)(2-x)与x轴所围图形的面积为( ）.

5.曲面 \(e ^ { z } - z + x y = 3\) 在点P(2,1,0)处的切平面方程为( ）.

6.设f(x)是周期为2的周期函数，且
\(f ( x ) { = } \left\{ \begin{array} { l l } { { 2 } , \ { - } 1 < x \le 0 } \\ { { x } ^ { 3 } , \quad 0 < x \le 1 , } \end{array} \right.\)
则f(x)的傅里叶级数在x=1处收敛于（ ）.

7.设f(u,v)为二元可微函数，若 \(z = f ( x ^ { y } , y ^ { x } )\) ，则
\(\frac { \partial z } { \partial x } = (\) ）.

8.设 \(L : y = x ^ { 2 } ( 0 \leq x \leq { \sqrt { 2 } } )\) ，则
\(\int _ { L } x d s = (\) ）.

\subsection{二、计算题（本大题共4小题，共30分）}\label{ux4e8cux8ba1ux7b97ux9898ux672cux5927ux9898ux51714ux5c0fux9898ux517130ux5206}

1．（本小题5分）求极限
\(\lim _ { x \to 0 ^ { + } } { \frac { \sqrt { x } } { 1 - e ^ { \sqrt { x } } } }\)

2．（本小题5分）计算积分
\(\int _ { 0 } ^ { \frac { \pi } { 2 } } { \frac { \cos x } { 1 + \sin ^ { 2 } x } } d x\)

3．（本小题10分）求幂级数
\(\sum _ { n = 1 } ^ { \infty } n x ^ { n - 1 }\) 的收敛域及和函数.

4．（本小题10分）计算曲线积分
\(\int _ { L } ( x ^ { 2 } + 3 y ) d x + ( y ^ { 2 } - x ) d y\)
,其中L为从点 \(\pmb { A } ( 4 , 0 )\) 到点0(0,0)的上半圆周
\(y = \sqrt { 4 x - x ^ { 2 } }\)

\subsection{三、综合题（本大题共3小题，每小题10分，共30分)}\label{ux4e09ux7efcux5408ux9898ux672cux5927ux9898ux51713ux5c0fux9898ux6bcfux5c0fux989810ux5206ux517130ux5206}

1.已知函数 \(f ( x , y ) = { \sqrt { | x y | } }\)

(1)求f(x,y)在(0,0)处的偏导数 \(f _ { x } ( 0 , 0 )\) 和
\({ f _ { y } } ( 0 , 0 )\)

(2)讨论f(x,y)在点(0,0)处的可微性.

2.讨论函数项级数
\(\sum _ { n = 1 } ^ { \infty } { \frac { x ^ { 2 0 2 1 } } { e ^ { n x } } }\)
在 \([ 0 , + \infty )\) 上的一致收敛性.

3.计算曲面积分
\(\int \int _ { \Sigma } { x z ^ { 2 } d y d z } + ( x ^ { 2 } y - z ^ { 2 } ) d z d x + ( 2 x y + y ^ { 2 } z ) d x d y\)
，其中 \(\Sigma\) 为上半球面\(z = \sqrt { 1 - x ^ { 2 } - y ^ { 2 } }\)
的上侧.

\subsection{四、证明题（本大题共5小题，每小题10分，共50分）}\label{ux56dbux8bc1ux660eux9898ux672cux5927ux9898ux51715ux5c0fux9898ux6bcfux5c0fux989810ux5206ux517150ux5206}

1.设
\(a > 0 , a _ { 1 } = \frac { 1 } { 2 } \left( a + \frac { 1 } { a } \right) , a _ { n + 1 } = \frac { 1 } { 2 } \left( a _ { n } + \frac { 1 } { a _ { n } } \right) , n = 1 , 2 , \cdots\)
．证明：数列 \(\{ a _ { n } \}\) 收敛,并且 其极限为1.

2.用一致连续的定义证明：函数 \(f ( x ) { = } \cos { \sqrt { x } }\) 在
\([ a , + \infty ) \ ( a > 0 )\) 上一致连续.

3.设f(x)在{[}a,b{]}上连续，在(a,b)内可微，且

\[
f ( a ) \cdot f ( b ) > 0 , f ( a ) \cdot f \left( \frac { a + b } { 2 } \right) < 0 .
\]

证明：存在一点 \(\xi \in ( a , b )\) ，使得
\(f ^ { \prime } ( \xi ) = f ( \xi )\)

4．证明：含参量反常积分
\(\int _ { 0 } ^ { + \infty } x e ^ { - \alpha x }\) dx在
\(\left[ \alpha _ { 0 } , + \infty \right) \left( \alpha _ { 0 } > 0 \right)\)
上一致收敛，但在\(( 0 , + \infty )\) 上非一致收敛.

5.设f(x)在 \([ a , b ] ( a < b )\) 上连续且单调递增.求证：

\[
\int _ { a } ^ { b } x f ( x ) d x \geq { \frac { a + b } { 2 } } \int _ { a } ^ { b } f ( x ) d x .
\]
\end{document}
